%% LyX 2.4.0~RC3 created this file.  For more info, see https://www.lyx.org/.
%% Do not edit unless you really know what you are doing.
\documentclass[french,titlepage]{extarticle}
\usepackage[T1]{fontenc}
\usepackage[utf8]{inputenc}
\usepackage{geometry}
\geometry{verbose,tmargin=2cm,bmargin=2cm,lmargin=2cm,rmargin=2cm}
\usepackage{mathrsfs}
\usepackage{amsmath}
\usepackage{amssymb}

\makeatletter
%%%%%%%%%%%%%%%%%%%%%%%%%%%%%% User specified LaTeX commands.
\usepackage{hyperref}   % pour les liens cliquables de la table des matières

%\usepackage[svgnames,dvipsnames]{xcolor}

\makeatletter
\@addtoreset{section}{part}   % reprendre à partir de 1 les sections des parties suivantes.
\makeatother

\usepackage{tocbibind} % pour afficher les références dans la table des matières

\makeatother

\usepackage{babel}
\begin{document}
\title{Mini-projet 5 -- Optimisation}
\author{Erhel CAYZAC, Lucie DUPOUY}
\date{2025 -- 2026}
\maketitle

\section{Étude du problème d'optimisation}
\begin{enumerate}
\item 
\item On connaît $p_{r}\in\mathbb{R}^{m_{r}}$, $f_{d}\in\mathbb{R}^{m_{d}}$.\\
Le problème peut donc se reformuler :\\
\begin{equation}
\min_{\begin{array}{c}
q\in\mathbb{R}^{l}\\
p,f\in\mathbb{R}^{m}\\
c_{eq}
\end{array}}\varphi(q,f,p)=\left\Vert Aq-f\right\Vert ^{2}+\left\Vert r\bullet q\bullet\vert q\vert+A^{T}p\right\Vert ^{2}
\end{equation}
\\
Avec pour contraintes :\\
\begin{equation}
\begin{cases}
Aq-f=0\\
r\bullet q\bullet\vert q\vert+A^{T}p=0\\
p=(p_{r},p_{d}) & p_{r}\text{ connu}\\
f=(f_{r},f_{d}) & f_{d}\text{ connu}
\end{cases}
\end{equation}
\\
\\
On a donc $n=3$ variables de décision, à savoir $q,f,p$.\\
\item On pose :\\
\[
\mathcal{L}:(q,\lambda)\in\mathbb{R}^{l}\times\mathbb{R}^{d}\longmapsto\frac{1}{3}q^{T}r\bullet q\bullet\vert q\vert+p_{r}^{T}A_{r}q+\lambda^{T}(A_{d}q-f_{d})
\]
\\
On cherche un point stationnaire de $\mathcal{L}$, ce qui donne les
contraintes suivantes :\\
\begin{equation}
\begin{cases}
\frac{\partial\mathcal{L}}{\partial q}=0\\
\frac{\partial\mathcal{L}}{\partial\lambda}=0
\end{cases}\Leftrightarrow\begin{cases}
r\bullet q\bullet\vert q\vert+A_{r}p_{r}^{T}+A_{d}^{T}\lambda=0\\
A_{d}q-f_{d}=0
\end{cases}
\end{equation}
\\
Les deux équations en jeux se retrouvent bien dans les équations (2)
et (4) de l'énoncé.\\
\item L'optimisation précédente permet d'obtenir $p_{r},q,f_{d}$. Pour
obtenir le vecteur $f$, on utilise la condition suivante :\\
\begin{equation}
Aq-f=0
\end{equation}
\\
Pour obtenir le vecteur $p$, on se sert de l'équation :\\
\begin{equation}
A^{T}p=-r\bullet q\bullet\vert q\vert
\end{equation}
\\
Cette équation devient :\\
\[
A_{d}^{T}p_{d}=-r\bullet q\bullet\vert q\vert-A_{r}^{T}p_{r}:=Q
\]
\\
On sait de plus que $A_{d}\in\mathscr{M}_{m_{d}l}(\mathbb{R})$. Il
s'agit d'une matrice de rang plein $m_{d}$, comme indiqué dans la
remarque 3. De plus, $l>m_{d}$. On extrait donc $m_{d}$ colonnes
linéairement indépendantes dans $A_{d}$. On note $\widetilde{A_{d}}$
la matrice inversible ainsi extraite de $A_{d}$, et $\widetilde{Q}$
la matrice des $m_{d}$ lignes extraites de $Q$ en question. On a
en particulier :\\
\[
\widetilde{A_{d}}^{T}p_{d}=\widetilde{Q}\Leftrightarrow p_{d}=\widetilde{A_{d}}^{-T}\widetilde{Q}
\]
\\
\\
\\
\end{enumerate}

\end{document}
